\begin{helper}
Let \(V\) be finite, let \(\nu\) be a measure on
\(\operatorname{Finset}(V)\times \mathbb R^V\), and fix an activation fiber
\(A\).  Assume the fixed-fiber cube law and lower-rectangle law used in
\noderef{RandomIndependentSetSampling}: the fiber \(A\) has all activated
priority coordinates in \([0,1]\) up to null error, and every lower rectangle
\(0\le \pi(v)\le t_v\), with \(t_v\in[0,1]\), has measure
\[
\nu[\omega_1=A]\,\operatorname{ofReal}\Bigl(\prod_{v\in V} t_v\Bigr).
\]
Let \(\gamma>0\), let \(a,b,c\in V\) be pairwise distinct, and let
\(S_x,S_y,S_z\subseteq V\) be pairwise disjoint subsets, each disjoint from
\(\{a,b,c\}\).  Then the measure of the event
\[
\omega_1=A,\qquad
\pi(a)\le \pi(b)\le \pi(c),
\]
together with
\[
\pi(w)<\pi(a)\ (w\in S_x),\qquad
\pi(w)<\pi(b)\ (w\in S_y),\qquad
\pi(w)<\pi(c)\ (w\in S_z),
\]
is at most
\[
\nu[\omega_1=A]\operatorname{ofReal}\left(
\frac{1}{\gamma^3}
\int_{0\le z\le y\le x\le\gamma}
\left(1-\frac{x}{\gamma}\right)^{|S_x|}
\left(1-\frac{y}{\gamma}\right)^{|S_y|}
\left(1-\frac{z}{\gamma}\right)^{|S_z|}\,dx\,dy\,dz\right).
\]
\end{helper}
\begin{proof}
Split according to the mass of the fixed activation fiber.  First suppose that
\(\nu[\omega_1=A]=\infty\).  Then the right hand side of the desired
inequality is infinite.  Indeed, on \(0\le z\le y\le x\le\gamma\) each factor
\(1-x/\gamma\), \(1-y/\gamma\), and \(1-z/\gamma\) lies in \([0,1]\), so the
integrand is nonnegative.  More than this, on the positive-volume subchamber
\[
0\le z\le y\le x\le \gamma/2
\]
each of these three factors is at least \(1/2\).  Since
\(|S_x|,|S_y|,|S_z|\) are finite, the integrand is bounded below on that
subchamber by the positive constant
\((1/2)^{|S_x|+|S_y|+|S_z|}\).  Therefore the displayed chamber integral is
strictly positive, and so is its product with \(1/\gamma^3\).  Multiplying the
positive extended nonnegative real obtained from this positive real by the
infinite fiber mass gives \(\infty\).  In this case the claimed upper bound is
therefore automatic.

It remains to consider the finite case
\(\nu[\omega_1=A]<\infty\).  This is exactly the finite fixed-fiber
product-cube calculation stated in
\noderef{SamplingTripleProductChamberIntegralFiniteFiber}.  The hypotheses of
that node are the present cube law, lower-rectangle law, finite fiber-mass
assumption, positivity of \(\gamma\), pairwise distinctness of \(a,b,c\), and
the disjointness assumptions on \(S_x,S_y,S_z\).  Applying that node gives the
same displayed chamber-integral upper bound in the finite case.  The two cases
therefore prove the aggregate estimate.
\end{proof}
